% !TeX root = ../../progress-sem1.tex
% 2.x ภาษาคอมพิวเตอร์ เครื่องมือ และเทคโนโลยี
\section{ภาษาคอมพิวเตอร์ เครื่องมือ และเทคโนโลยี}

\subsection{แบบจำลองภาษาขนาดใหญ่ (Large Language Models: LLM)}

\thaijustify{
  แบบจำลองภาษาขนาดใหญ่ (Large Language Models: LLM) คือ เทคโนโลยีปัญญาประดิษฐ์ที่ฝึกด้วยข้อมูลข้อความจำนวนมาก ทำให้สามารถเข้าใจภาษาธรรมชาติ วิเคราะห์บริบท ให้เหตุผล และสร้างเนื้อหาหรือโค้ดได้ จึงสามารถนำมาช่วยลดภาระงานและสนับสนุนกระบวนการพัฒนาระบบได้อย่างมีประสิทธิภาพ
}

\thaijustify{
  โครงงานนี้ใช้ LLM ผ่าน CPE AI Gateway \cite{cpe_ai_gateway} สำหรับการสร้าง UI และประมวลผลภาษาธรรมชาติ โดยออกแบบกระบวนการแบบ AI Orchestration เพื่อเลือกโมเดลให้เหมาะสมกับแต่ละงาน โมเดลที่นำมาใช้ ได้แก่ muse-glimmer-30b และ qwen3.8-27b พร้อมเปรียบเทียบกับกลุ่ม Frontier Model เพื่อพิจารณาจุดเด่นและความเหมาะสมในการใช้งาน
}

% \thaijustify{
%   muse-glimmer-30b เป็นโมเดล open-weight ขนาดประมาณ 30 พันล้านพารามิเตอร์ที่มีความสมดุลระหว่างความสามารถในการให้เหตุผลและความเร็วในการประมวลผล โดยเฉพาะการสร้างผลลัพธ์แบบมีโครงสร้าง จึงเหมาะกับงาน Prompt-to-Design ที่แปลงคำสั่งภาษาธรรมชาติเป็นโครงสร้าง UI แบบสองทางเลือก (two-alternative UI) อีกทั้งยังใช้ทรัพยากรน้อยกว่าการใช้ Frontier Model ในทุกขั้นตอน \cite{meta2026museglimmer}
% }

\thaijustify{
   muse-glimmer-30b (Muse Glimmer) เป็นโมเดล open-weight ขนาดประมาณ 30 พันล้านพารามิเตอร์จาก Meta Superintelligence Labs ออกแบบมาสำหรับงานเอเจนต์ที่ทำงานต่อเนื่องบนเครื่องภายใน และรันบนฮาร์ดแวร์ระดับผู้บริโภคได้~\cite{meta2026museglimmer}จึงใช้ทรัพยากรน้อยกว่าการเรียกใช้ Frontier Model ในทุกขั้นตอน โครงงานนี้จึงเลือกใช้โมเดลนี้สำหรับงาน Prompt-to-Design ซึ่งแปลงคำสั่งภาษาธรรมชาติเป็นโครงสร้าง UI แบบสองทางเลือก (two-alternative UI)
}

% \thaijustify{
%   qwen3.8-27b เป็นโมเดลขนาด 27 พันล้านพารามิเตอร์ที่เด่นด้านการเข้าใจและประมวลผลภาษาธรรมชาติ จึงเหมาะกับการวิเคราะห์ข้อเสนอแนะจากผู้ทดสอบที่มักเป็นภาษาพูดและไม่มีโครงสร้าง ระบบสามารถใช้โมเดลนี้เพื่อสรุป จัดกลุ่ม และสกัดข้อมูลเชิงลึกจากข้อเสนอแนะเพื่อนำไปสร้าง Persona โดยใช้ทรัพยากรน้อยกว่าโมเดลขนาดใหญ่ \cite{qwentheam2026qwen38}
% }

\thaijustify{
  qwen3.8-27b (Qwen3.8-27B) เป็นโมเดล dense แบบ open-weight ขนาด 27 พันล้านพารามิเตอร์จาก Alibaba เผยแพร่ภายใต้สัญญาอนุญาต Apache~2.0 รับข้อมูลได้ทั้งข้อความ ภาพ และวิดีโอ มี context window 262,144 tokens ซึ่งขยายได้ถึงประมาณ 1 ล้าน tokens~\cite{qwentheam2026qwen38} context ที่ยาวนี้ทำให้ส่งข้อเสนอแนะจำนวนมากจากผู้เข้าร่วมทดสอบหลายคนเข้าไปวิเคราะห์ในคำขอเดียวได้ โครงงานนี้จึงใช้โมเดลนี้สรุป จัดกลุ่ม และสกัดข้อมูลเชิงลึกจากข้อเสนอแนะเพื่อนำไปสร้าง Persona Artifact ทั้งนี้ ผลการประเมินที่เผยแพร่ในช่วงเปิดตัวเป็นตัวเลขที่ผู้พัฒนาโมเดลรายงานเอง จึงต้องประเมินความสามารถของโมเดลกับข้อเสนอแนะภาษาไทย ซึ่งมักเป็นภาษาพูดและไม่มีโครงสร้าง
}

\thaijustify{
  และกลุ่ม Frontier Model มีตัวอย่างในช่วงที่จัดทำโครงงานนี้ ได้แก่ GPT-6 Astra ของ OpenAI ซึ่งเน้นความสามารถด้านการใช้งานคอมพิวเตอร์และวิศวกรรมซอฟต์แวร์ \cite{openai2026gpt6astra} Claude Opus 5.5 ของ Anthropic ซึ่งเน้นงานเขียนโค้ดแบบเอเจนต์ (Agentic Coding) และงานเชิงความรู้ \cite{anthropic2026opus55} และ Gemini 3.8 Flash ของ Google ซึ่งพัฒนาด้านวิศวกรรมซอฟต์แวร์และงานแบบเอเจนต์ โดยยังคงความเร็วสูงและต้นทุนต่ำ \cite{google2026gemini38flash} โมเดลกลุ่มนี้ให้คุณภาพผลลัพธ์สูงในงานที่ซับซ้อน แต่ส่วนใหญ่เป็นโมเดลแบบปิด (Closed-weight) ที่เข้าถึงได้ผ่าน API เชิงพาณิชย์เท่านั้น มีค่าใช้จ่ายตามจำนวนโทเคน และต้องส่งข้อมูลออกไปประมวลผลบนระบบภายนอก ซึ่งเป็นข้อพิจารณาสำคัญเมื่อข้อมูลที่ต้องประมวลผลคือข้อเสนอแนะของผู้เข้าร่วมทดสอบที่อยู่ภายใต้กฎหมายคุ้มครองข้อมูลส่วนบุคคล (PDPA) ด้วยเหตุนี้ โครงงานจึงใช้ Frontier Model เป็นเกณฑ์อ้างอิง (Baseline) ในการประเมินคุณภาพผลลัพธ์ของ muse-glimmer-30b และ qwen3.8-27b แทนการใช้เป็นโมเดลหลักในทุกขั้นตอน
}

\begin{table}[H]
  \centering
  \small
  \caption{เปรียบเทียบความสามารถและจุดด้อยของแบบจำลองภาษาขนาดใหญ่แต่ละรุ่น}
  \label{tab:ai-models-comparison}
  \begin{tabularx}{\textwidth}{
      >{\raggedright\arraybackslash}p{0.20\textwidth}
      |
      >{\raggedright\arraybackslash}X
      >{\raggedright\arraybackslash}X
      >{\raggedright\arraybackslash}X
    }
    \hhline{=|===}
    \textbf{Model} &
    \textbf{ประเภทงาน} &
    \textbf{จุดเด่น} &
    \textbf{จุดด้อย} \\
    \hhline{=|===}
    
    \texttt{muse-glimmer-30b} \newline \textit{(Context: 131K tokens)} &
    Prompt-to-Design และงานที่ต้องใช้การให้เหตุผล &
    มีความสามารถในการทำตามคำสั่งแบบมีโครงสร้าง (Structured output) ได้อย่างแม่นยำ ใช้ทรัพยากรน้อย และตอบสนองรวดเร็ว &
    ความสามารถในการวิเคราะห์ข้อมูลที่ซับซ้อนมาก ๆ หรือความรู้รอบกว้าง ยังเป็นรองโมเดลขนาดใหญ่ (Frontier Models) และมี Context Window ที่จำกัด \\

    \hline
    
    \texttt{qwen3.8-27b} \newline \textit{(Context: 256K - 1M tokens)} &
    การวิเคราะห์ข้อความ การสนทนา และการประมวลผลภาษาธรรมชาติ &
    เข้าใจภาษาพูดและบริบทที่คลุมเครือได้ดีเยี่ยม สรุปใจความ จัดกลุ่มข้อเสนอแนะ และดึงข้อมูลเชิงลึกได้อย่างเป็นธรรมชาติ &
    ขีดความสามารถด้านการคำนวณเชิงตรรกะซับซ้อน หรือการเขียนโค้ดขั้นสูง ยังไม่เทียบเท่ากลุ่ม Frontier Model \\

    \hline

    \texttt{Claude Opus 5.5} \newline \textit{(Frontier Model)} \newline \textit{(Context: 1M tokens)} &
    งานวิเคราะห์เชิงลึกขั้นสูง การเขียนและตรวจสอบโค้ดโปรแกรม &
    มีระบบ Adaptive thinking ที่ช่วยให้เหตุผลได้อย่างละเอียดลึกซึ้ง การใช้งานเครื่องมือ (Tools) มีความเสถียรและแม่นยำสูง &
    มีต้นทุนการเรียกใช้งาน (Cost) ที่ค่อนข้างสูง และระยะเวลาประมวลผล (Latency) อาจใช้เวลานานกว่าโมเดลทั่วไป \\

    \hline
    
    \texttt{Gemini 3.8 Flash} \newline \textit{(Frontier Model)} \newline \textit{(Context: 1M tokens)} &
    งานประมวลผลแบบ Multimodal และงานที่ต้องการผลลัพธ์แบบเรียลไทม์ &
    ประมวลผลได้อย่างรวดเร็วมาก (Fastest output) รองรับการรับข้อมูลทั้งข้อความ ภาพ เสียง และวิดีโอได้ดีเยี่ยมในราคาประหยัด &
    ประสิทธิภาพในการให้เหตุผลเชิงตรรกะสำหรับงานที่ซับซ้อนมาก ๆ อาจยังเป็นรองโมเดลในระดับ Pro หรือระดับที่สูงกว่า \\

    \hline
    
    \texttt{GPT-6 Astra} \newline \textit{(Frontier Model)} \newline \textit{(Context: 1M - 1.1M tokens)} &
    การประสานงานแบบ Multi-agent และการแก้ปัญหาที่ซับซ้อนขั้นสุด &
    มีศักยภาพการให้เหตุผลระดับสูงสุด สามารถควบคุมระบบคอมพิวเตอร์และทำงานซ้อนกันหลายขั้นตอน (Complex reasoning) ได้ดีเยี่ยม &
    มีต้นทุนการใช้งานที่สูงที่สุดในกลุ่ม อาจเกินความจำเป็นสำหรับระบบทั่วไป และมักมีข้อจำกัดด้านโควตาการใช้งานอย่างเข้มงวด \\

    \hline
  \end{tabularx}
\end{table}

% \begin{table}[H]
%   \centering
%   \small
%   \caption{เปรียบเทียบจุดเด่นและจุดด้อยของแบบจำลองภาษาขนาดใหญ่แต่ละรุ่น}
%   \label{tab:ai-models-comparison}
%   \begin{tabularx}{\textwidth}{
%       >{\raggedright\arraybackslash}p{0.20\textwidth}
%       |
%       >{\raggedright\arraybackslash}X
%       >{\raggedright\arraybackslash}X
%       >{\raggedright\arraybackslash}X
%     }
%     \hhline{=|===}
%     \textbf{Model} &
%     \textbf{ประเภทงาน} &
%     \textbf{จุดเด่น} &
%     \textbf{จุดด้อย} \\
%     \hhline{=|===}
 
%     % TODO: ยืนยัน context จาก https://huggingface.co/meta-models/Muse-Glimmer-30B
%     \texttt{muse-glimmer-30b} \newline \textit{(Context: 131K tokens)} &
%     Prompt-to-Design (สร้างและปรับปรุง two-alternative UI) &
%     โมเดล open-weight ขนาด 30B ที่ออกแบบมาสำหรับงานเอเจนต์ รันบนฮาร์ดแวร์ระดับผู้บริโภคได้ จึงใช้ทรัพยากรน้อย~\cite{meta2026museglimmer} &
%     Context window สั้นกว่ากลุ่ม Frontier Model และความแม่นยำในการสร้างผลลัพธ์แบบมีโครงสร้างสำหรับงาน Prompt-to-Design ยังต้องยืนยันด้วยชุดข้อมูลประเมินของโครงงาน \\
 
%     \hline
 
%     \texttt{qwen3.8-27b} \newline \textit{(Context: 262K ขยายได้ถึง 1M tokens)} &
%     การวิเคราะห์ สรุป และจัดกลุ่มข้อเสนอแนะ และการสังเคราะห์ข้อมูลสำหรับ Persona &
%     โมเดล dense แบบ open-weight (Apache~2.0) รับข้อมูลได้ทั้งข้อความ ภาพ และวิดีโอ มี context ยาวพอสำหรับรวมข้อเสนอแนะจำนวนมากในคำขอเดียว~\cite{qwentheam2026qwen38} &
%     ผลประเมินที่เผยแพร่ช่วงเปิดตัวเป็นตัวเลขที่ผู้พัฒนารายงานเอง และความสามารถกับข้อเสนอแนะภาษาไทยแบบภาษาพูดยังต้องประเมินด้วยชุดข้อมูลของโครงงาน \\
 
%     \hline
 
%     % TODO: ยืนยัน context ที่ https://platform.claude.com/docs/en/models/opus-5-5/overview
%     \texttt{Claude Opus 5.5} \newline \textit{(Frontier Model)} \newline \textit{(Context: 1M tokens)} &
%     งานวิเคราะห์เชิงลึก การเขียนโค้ดแบบเอเจนต์ และงานเชิงความรู้ &
%     ผู้พัฒนาระบุว่าเป็นรุ่น Opus ที่ทำงานด้านโค้ด เอเจนต์ และงานเชิงความรู้ได้ดีที่สุด ต้นทุนต่องานต่ำกว่า Opus~5 ราว 40\% และสร้างผลลัพธ์เร็วขึ้นกว่า 30\%~\cite{anthropic2026opus55} &
%     เป็นโมเดลแบบปิด (Closed-weight) ต้องส่งข้อมูลออกไปประมวลผลภายนอก และมีราคา 4/20 ดอลลาร์สหรัฐต่อล้าน tokens (input/output) ซึ่งสูงกว่าโมเดล open-weight มาก \\
 
%     \hline
 
%     % TODO: ยืนยัน context ของรุ่น Flash (หน้า docs ที่พบระบุ 1,048,576 tokens สำหรับรุ่น Flash Cyber)
%     \texttt{Gemini 3.8 Flash} \newline \textit{(Frontier Model)} \newline \textit{(Context: 1M tokens)} &
%     งานวิศวกรรมซอฟต์แวร์ งานแบบเอเจนต์ และงานที่ต้องการผลลัพธ์รวดเร็ว &
%     พัฒนาด้านวิศวกรรมซอฟต์แวร์และงานแบบเอเจนต์ โดยคงความเร็วสูงและต้นทุนต่ำ~\cite{google2026gemini38flash} &
%     % TODO: ช่องนี้ยังไม่มีแหล่งรองรับ
%     ประสิทธิภาพในการให้เหตุผลสำหรับงานที่ซับซ้อนมาก ๆ อาจยังเป็นรองโมเดลระดับ Pro หรือสูงกว่า \\
 
%     \hline
 
%     % TODO: ยืนยัน context จากหน้า model ของ OpenAI API
%     \texttt{GPT-6 Astra} \newline \textit{(Frontier Model)} \newline \textit{(Context: 1M - 1.1M tokens)} &
%     งานที่ใช้คอมพิวเตอร์ (Computer use) วิศวกรรมซอฟต์แวร์ และงานวิชาชีพที่ซับซ้อน &
%     ผู้พัฒนาระบุว่าทำผลงานระดับ state-of-the-art ด้าน computer use, browsing, วิศวกรรมซอฟต์แวร์ และงานวิชาชีพ~\cite{openai2026gpt6astra} &
%     % TODO: ยืนยันราคาจากหน้า pricing ทางการของ OpenAI
%     มีต้นทุนการใช้งานสูงที่สุดในกลุ่ม (10/50 ดอลลาร์สหรัฐต่อล้าน tokens) และอาจเกินความจำเป็นสำหรับระบบทั่วไป \\
 
%     \hline
%   \end{tabularx}
% \end{table}

\thaijustify{
  ในกระบวนการทำงาน AI Orchestration จะเลือกโมเดลตามประเภทของงาน เช่น คำขอที่เกี่ยวข้องกับการสร้างหรือปรับปรุง UI จากความต้องการทางการออกแบบจะถูกส่งไปยัง muse-glimmer-30b ขณะที่งานสนทนา การสรุปและการวิเคราะห์ข้อเสนอแนะ รวมถึงการสังเคราะห์ข้อมูลสำหรับ Persona Artifact จะใช้ qwen3.8-27b เป็นหลัก
}
\thaijustify{
  การแบ่งหน้าที่ดังกล่าวช่วยให้แต่ละโมเดลถูกใช้งานตามความสามารถที่เหมาะสมและทำให้กระบวนการประมวลผลของระบบมีโครงสร้างชัดเจนมากกว่าการใช้โมเดลเดียวสำหรับงานทุกประเภท
}

\subsection{เทคโนโลยีการสร้างส่วนต่อประสานผู้ใช้โดยกำหนด Design System}
\thaijustify{
  งานวิจัยที่ศึกษาการให้ Generative AI สร้าง UI ให้ตรงกับ design system มีการศึกษาเชิงเปรียบเทียบ context engineering สามแบบสำหรับ LLM agent ได้แก่ instruction-based, context-based และ registry-based \cite{designsystem_compliant_ui_2026} ผลพบว่าการให้ agent ใช้ component สำเร็จรูปจาก registry ให้ความสอดคล้องกับ design system สูงและขยายขนาดได้ดีกว่า เพราะเปลี่ยนจากการให้ LLM ตีความข้อความไปเป็นการตรวจสอบระดับ component และมีงานอีกกลุ่มหนึ่งใช้ข้อจำกัดด้านการออกแบบเป็นตัวกำหนด เช่น UIDEC~\cite{uidec_2025} ที่ให้นักออกแบบระบุรายละเอียดโครงการ ได้แก่ วัตถุประสงค์ กลุ่มเป้าหมาย อุตสาหกรรม และสไตล์การออกแบบ แล้วใช้ Generative AI สร้างตัวอย่างการออกแบบที่หลากหลายภายใต้ข้อจำกัดนั้นโดยนักออกแบบแทบไม่ต้องเขียน prompt เอง และ PrototypeFlow~\cite{prototypeflow_humanai_ui_design} ที่สร้าง theme ที่แก้ไขได้และดึงความรู้ด้านการออกแบบมาช่วยให้สี สไตล์ และ layout มีความสอดคล้องกัน จากบทความเหล่านี้ โครงงานนี้จึงนำแนวคิดนี้มาประยุกต์ใช้ในการสร้าง component และ material ที่ทำหน้าที่เป็นโครงสร้างพื้นฐานให้ LLM หรือนักพัฒนาสามารถนำไปหยิบใช้และทำงานต่อได้ทันทีโดยไม่ต้องเสียเวลาเขียนโค้ดหรือจัดการรายละเอียดด้วยตนเองอีกครั้ง
}

% /2.1.6 - 2.2.3 เอามารวมกันเลย
